This thesis presents a layout-aware discriminative approach to automatic
key-value pair (KVP) extraction from business documents, evaluated on the
KVP10k benchmark (ICDAR~2024, Naparstek et al.).

The KVP10k paper argued, without direct empirical evidence, that
encoder-based models such as LayoutLMv3 are insufficient for KVP linking,
and consequently evaluated only a generative LMDX-style baseline.
This thesis empirically tests that claim by extending LayoutLMv3
(125M parameters) with a lightweight Biaffine Linker module and a two-phase
training strategy: entity pre-training followed by joint fine-tuning.

We evaluate two linking granularities. Token-level linking (V1) yields
near-zero link F1 (0.008) due to extreme class imbalance ($\sim$450:1
negative-to-positive ratio). Span-level linking (V2) reduces this to
$\sim$5:1 but initially produces link F1 of only 0.007 due to three
compounding data-pipeline bugs identified through systematic diagnostic
experiments. After correcting these bugs, the final V4 model achieves
entity F1~$=0.827$ and, under the released KVP10k macro benchmark, regular-KVP
F1~$=0.334$ for text and $0.309$ for text+location. The corresponding
published scores are $0.659$ and $0.611$. In its primary configuration V4 emits
regular linked pairs only; a post-hoc decode-time recovery optionally extends
coverage to the benchmark's unkeyed and unvalued categories at markedly lower
accuracy.
Diagnostic analyses localize the remaining regular-pair gap primarily to
relation modeling, while the layout-grounded architecture prevents generation
of text absent from the OCR input.

\textbf{Keywords:} Document Understanding, Key-Value Pair Extraction,
LayoutLMv3, Span-Level Linking, Biaffine Linker, Multimodal Learning,
Information Extraction
